% part: intuitionistic-logic
% chapter: propositions-as-types
% section: introduction

\documentclass[../../../include/open-logic-section]{subfiles}

\begin{document}

\olfileid{int}{pty}{int}

\olsection{ভূমিকা}

ল্যাম্বডা ক্যালকুলাস গণনার একটি সরল মডেল। এতে চলরাশি দিয়ে গড়া
পদের উপর কাজ হয়। $N$ ও $M$ পদ হলে $(NM)$-ও পদ
(“$N$-কে~$M$-এ প্রয়োগ”)। $N$ পদ হলে
$\lambd[x][N]$-ও পদ। $N$-তে চলরাশি~$x$ থাকলে
$\lambd[x][N]$-এর মধ্যে তার সংশ্লিষ্ট $x$-এর
আবির্ভাবগুলি $\lambd[x]$ অপারেটর দ্বারা বদ্ধ হয়;
তাই আর মুক্ত থাকে না। $(\lambd[x][N]M)$ আকারের
পদকে $\beta$-হ্রাসে $\Subst{N}{M}{x}$ পাওয়া যায়
(অর্থাৎ~$N$-এ $x$-এর প্রতিটি মুক্ত আবির্ভাবের
জায়গায় $M$ বসে)। $N$-এর কোনো উপপদে
$\beta$-হ্রাস করে $N'$ পাওয়া গেলে $N$
পদটি $N'$-এ $\beta$-রূপান্তরিত হয়; লিখি
$N \redone N'$। $\lambd$-ক্যালকুলাসে গণনা হলো
$N \redone N_1 \redone N_2 \redone \dots$
আকারের অনুক্রম। অনুক্রমটি এমন $N_k$ পদে
শেষ হলে, যার কোনো উপপদে $\beta$-হ্রাস আর সম্ভব নয়,
তাকে~$N$-এর \emph{স্বাভাবিক রূপ} বলে। $\lambd$-ক্যালকুলাসে
গণনাকে এমন অনুক্রম হিসেবেই ভাবি। স্বাভাবিক রূপ
পাওয়া গেলে গণনা থামে, এবং সেই স্বাভাবিক রূপই
গণনার ফল। গণনার এই সরল মডেলটি টুরিং-সম্পূর্ণ:
প্রতিটি আংশিক পুনরাবৃত্ত অপেক্ষক $\lambd$-ক্যালকুলাসের
কোনো পদ দিয়ে গণনা করা যায়।

হাস্কেল কারি ও উইলিয়াম হাওয়ার্ড স্বাধীনভাবে
স্বাভাবিক নিষ্পাদন ও $\lambd$-ক্যালকুলাসের
ঘনিষ্ঠ সাদৃশ্য খুঁজে পান। স্বজ্ঞায়
$\lambd[x][N]$ পদটি একটি অপেক্ষক: $x$
তার যুক্তি; $x$ দেওয়া থাকলে মান কীভাবে গণনা
হবে, তা~$N$ জানায়। $(\lambd[x][N]M)$
পদে অপেক্ষক $\lambd[x][N]$-কে যুক্তি~$M$-এ
প্রয়োগ করা হয়। $\beta$-হ্রাসে তার মান
নির্ণয়ের রীতি: $N$-এ $x$-এর জায়গায়~$M$
বসাই (তারপর পাওয়া পদটির মূল্যায়ন চালিয়ে যাই)।
এবার অপেক্ষকগুলির নির্দিষ্ট সংজ্ঞাক্ষেত্র ও
মানক্ষেত্র ভাবি। $\lambd[x][N]$-এর সংজ্ঞাক্ষেত্র
$A$ এবং মানক্ষেত্র~$B$ হলে চলরাশি~$x$
শুধু~$A$ থেকে যুক্তি নেয়, আর $N$ শুধু~$B$-তে
মান দেয়। ফলে $\lambd[x][N]$ যদি $A \to B$
অপেক্ষক হয় এবং $M \in A$ হয়, তবেই
$(\lambd[x][N]M)$ অর্থবহ। এই তথ্য
$\lambd$-ক্যালকুলাসে যোগ করলে পাই
\emph{টাইপযুক্ত} $\lambd$-ক্যালকুলাস।
টাইপযুক্ত $\lambd$-ক্যালকুলাসে সব
$(\lambd[x][N]M)$ রাশি বৈধ নয়;
শুধু $\lambd[x][N]$ ও $M$-এর “টাইপ”
মিললে বৈধ—অর্থাৎ $\lambd[x][N]$-এর টাইপ
$A \lif B$ এবং $M$-এর টাইপ~$A$।
$\lambd[x][N]$-এর টাইপ $A \to B$ হওয়ার জন্য
$x$-এর টাইপ $A$ ও $N$-এর টাইপ~$B$ হতে হবে।
সাধারণভাবে সব $(M_1M_2)$ রাশি বৈধ নয়।
যে $(M_1M_2)$ পদে $M_1$-এর টাইপ
$A \to B$ এবং $M_2$-এর টাইপ~$A$,
সেটিই বৈধ; তখন $(M_1M_2)$-র টাইপ~$B$।
% BN-SRC-893: translation-local only-if direction restored.

এই টাইপ-নির্ধারণ বিধিগুলির সঙ্গে স্বাভাবিক নিষ্পাদনের
!!{derivation}s-এর স্পষ্ট অনুরূপতা আছে: টাইপের
জায়গায় !!{formula}s, আর !!{derivation}s-এর
জায়গায় $\lambd$-পদ। যেমন,
$!A \lif !B$-র !!a{derivation}-এর অনুরূপ
অপেক্ষক-টাইপ $!A \lif !B$-যুক্ত
$\lambd[\typeof{x}{!A}][N]$ পদ।
স্বাভাবিক নিষ্পাদনের অনুমানবিধি টাইপযুক্ত
ল্যাম্বডা ক্যালকুলাসের টাইপ-বিধির অনুরূপ; যেমন,
\begin{prooftree}
  \AxiomC{$\Discharge{!A}{x}$}
  \DeduceC{$!B$}
  \DischargeRule{\Intro\lif}{x}
  \UnaryInfC{$!A \lif !B$}
  \DisplayProof\bottomAlignProof
  \quad\text{অনুরূপ}\quad
  \Axiom$x: !A \fCenter N: !B$
  \UnaryInf$ \fCenter \lambd[\typeof{x}{!A}][N] : !A \lif !B$
\end{prooftree}
ডান দিকের বিধির অর্থ: $x$-এর টাইপ~$!A$ এবং
$N$-এর টাইপ~$!B$ হলে
$\lambd[\typeof{x}{!A}][N]$-এর টাইপ
$!A \lif !B$।

$\Elim\lif$ বিধির অনুরূপ হলো যৌগিক
পদের টাইপ-নির্ধারণ বিধি:
\[
  \AxiomC{$!A \lif !B$}
  \AxiomC{$!A$}
  \RightLabel{\Elim\lif}
  \BinaryInfC{$!B$}
  \DisplayProof
  \quad\text{অনুরূপ}\quad
  \Axiom$\fCenter M_1: !A \lif !B$
  \Axiom$\fCenter M_2: !A$
  \BinaryInf$ \fCenter (M_1M_2) : !B$
  \DisplayProof
\]
\Intro\lif{} বিধির ঠিক পরেই
\Elim\lif{} বিধি এলে !!{derivation} সরল করা যায়:
\begin{gather*}
  \AxiomC{$\Discharge{!A}{x}$}
  \DeduceC{$!B$}
  \DischargeRule{\Intro\lif}{x}
  \UnaryInfC{$!A \lif !B$}
  \AxiomC{}
  \DeduceC{$!A$}
  \RightLabel{\Elim\lif}
  \BinaryInfC{$!B$}
  \DisplayProof
  \quad\redone\quad
  \AxiomC{}
  \DeduceC{$!A$}
  \DeduceC{$!B$}
  \DisplayProof
\end{gather*}
এর অনুরূপ ল্যাম্বডা পদে $\beta$-হ্রাস:
\[
(\lambd[\typeof{x}{!A}][N]M) \quad\redone\quad
\Subst{N}{M}{x}.
\]

$\land$-এর বিধি ও “গুণন” টাইপের মধ্যে এবং
$\lor$-এর বিধি ও “সমষ্টি” টাইপের মধ্যেও
একই ধরনের অনুরূপতা আছে।

সরল টাইপযুক্ত ল্যাম্বডা ক্যালকুলাসের পদ ও
স্বাভাবিক নিষ্পাদনের !!{derivation}s-এর এই
অনুরূপতার নাম “কারি--হাওয়ার্ড”, “প্রমাণ-হিসেবে-প্রোগ্রাম”
বা “সূত্র-হিসেবে-টাইপ” অনুরূপতা। !!{formula}s
(প্রস্তাবনা)-র সঙ্গে টাইপ এবং প্রমাণের সঙ্গে
পদের অনুরূপতা ছাড়াও সম্পূর্ণ তালিকাটি এমন:
\begin{center}
  \begin{tabular}{c c c}
    যুক্তি & প্রোগ্রাম \\
    \hline
    প্রস্তাবনা/!!{formula} & টাইপ \\
    প্রমাণ & পদ \\
    অনুমান & চলরাশি \\
    অবমুক্ত অনুমান & বদ্ধ চলরাশি \\
    অবমুক্ত নয় এমন অনুমান & মুক্ত চলরাশি \\
    $!A \lif !B$ & অপেক্ষক-টাইপ \\
    $!A \land !B$ & গুণন-টাইপ \\
    $!A \lor !B$ & সমষ্টি-টাইপ \\
    $\lfalse$ & ফাঁকা টাইপ \\
    \hline
  \end{tabular}
\end{center}
% BN-NORM-176: translated display/caption text retained; optional annotation withdrawn.

কারি--হাওয়ার্ড অনুরূপতা স্বয়ংক্রিয় প্রমাণ-সহায়ক ও
প্রোগ্রামের টাইপ-পরীক্ষকের অন্যতম ভিত্তি। কারণ,
কোনো প্রস্তাবনার সাক্ষী প্রমাণ যাচাই করা (যেমন
উপরে করেছি) মানে একটি প্রোগ্রাম (পদ)-এর ঘোষিত
টাইপ আছে কি না যাচাই করা।
\end{document}
